\chapter{Derivación en sentido débil}
\noindent
Como motivacion para estudiar la derivada débil, si tenemos la ecuación de ondas:
\begin{equation*}
    \left\{\begin{array}{ll}
            u_{tt} - c^2 u_{xx} = 0 , &\text{en\ } \left[0,\infty\right[\times \mathbb{R} \\
                u(t=0,x) = \varphi(x), &\text{en\ } \mathbb{R} \\
                u_t(t=0,x) = 0 , & \text{en\ } \mathbb{R}
    \end{array}\right.
\end{equation*}
La solución se escribe así:
\begin{equation*}
    u(t,x) = \frac{\varphi(x+ct)+\varphi(x-ct)}{2}
\end{equation*}
Y esta puede presentar ``picos'' en su gráfica, es decir, puntos aislados en los que no es derivable.

\section{Definición de derivada débil}
\noindent
Sea $\left]a,b\right[$ un intervalo acotado, sabemos que para $f\in C^1(\left]a,b\right[)$ tenemos:
\begin{equation}\label{eq:int_partes}
    \int_{a}^{b} f'(x)\varphi(x)~dx = -\int_{a}^{b} f(x)\varphi'(x)~dx 
\end{equation}
Para $\varphi \in C^1_c\left(\left]a,b\right[\right)$, luego $\varphi(a) = 0 = \varphi(b)$, ya que $\sop \varphi \subsetneq \left]a,b\right[$ por ser $\sop\varphi$ compacto\footnote{Recordamos que $\sop\varphi = \overline{\{x\in \left]a,b\right[ : \varphi(x)\neq 0\}}$.}. También es claro que $\sop\varphi' \subset \sop\varphi$.\\

\noindent
Para que la fórmula~\eqref{eq:int_partes} tenga sentido, necesitamos exigir:
\begin{equation*}
    f'\in L^1(\sop\varphi), \qquad f\in L^1(\sop\varphi')
\end{equation*}
Así, bastaría con $f,f' \in L^1_{\text{loc}}(\left]a,b\right[)$, pues tanto $\sop\varphi$ como $\sop\varphi'$ son compactos.

\begin{definicion}
    Sea $f\in L^1_\text{loc}(\left]a,b\right[)$, diremos que admite derivada débil si existe $g\in L^1_\text{loc}(\left]a,b\right[)$ tal que:
    \begin{equation*}
        \int_{a}^{b} f(x)\varphi'(x)~dx  = -\int_{a}^{b} g(x)\varphi(x)~dx  \qquad \forall \varphi \in C_c^1(\left]a,b\right[)
    \end{equation*}
    En cuyo caso, diremos que $g$ es la derivada débil de $f$, notado por $g = f'$.
\end{definicion}

\noindent
En primer lugar tenemos que ver que la derivada débil está bien definida, es decir, que supuesto que $f\in L^1_\text{loc}(\left]a,b\right[)$ tiene derivada débil, esta ha de ser única, salvo un conjunto de medida nula:
\begin{proof}
    Para $f\in L^1_\text{loc}(\left]a,b\right[)$, por reducción al absurdo, suponemos que tenemos $g_1,g_2\in L^1_\text{loc}(\left]a,b\right[)$ tales que:
    \begin{multline*}
        \int_{a}^{b} f(x)\varphi'(x)~dx  = -\int_{a}^{b} g_1(x)\varphi(x)~dx, \quad 
        \int_{a}^{b} f(x)\varphi'(x)~dx  = -\int_{a}^{b} g_2(x)\varphi(x)~dx  \\ \forall \varphi \in C_c^1(\left]a,b\right[) 
    \end{multline*}
    De aquí deducimos que:
    \begin{equation*}
        \int_{a}^{b} (g_1(x)-g_2(x))\varphi(x)~dx = 0 \qquad \varphi \in C_c^1(\left]a,b\right[)
    \end{equation*}
    Por una versión del Lema Fundamental del Cálculo de Variaciones concluimos que $g_1-g_2\equiv 0$ casi por doquier, por lo que $g_1 = g_2$ casi por doquier.
\end{proof}

\begin{observacion}
    Observamos que si $f\in C^1(\left]a,b\right[)$, entonces $f$ es débilmente derivable, y su derivada débil coincide con su derivada habitual, salvo un conjunto de medida nula.\\

    \noindent
    Normalmente diremos que la derivada débil coincide con la habitual, puesto que lo que sí tenemos es que la derivada débil pertenece a la misma clase de equivalencia que la derivada usual en $L^1_\text{loc}(\left]a,b\right[)$, y de esa clase podemos tomar la derivada usual como representante canónico de cualquier elemento de la clase. Así, a partir de ahora omitiremos el ``salvo un conjunto de medida nula'', entiendo que trabajamos con funciones en $L^1_\text{loc}(\left]a,b\right[)$.
\end{observacion}

\begin{prop}
    Si $f_1,f_2\in L^1_\text{loc}(\left]a,b\right[)$ son débilmente derivables y $\lm_1,\lm_2\in \mathbb{R}$, entonces:
    \begin{equation*}
        (\lm_1f_2 + \lm_2f_2)' = \lm_1f_1' + \lm_2f_2'
    \end{equation*}
    \begin{proof}
        Se deduce directamente de la linealidad de la integral, puesto que la derivada débil se define a partir de una función que verifica una igualdad integral.
    \end{proof}
\end{prop}

\begin{observacion}
    La derivada en sentido débil no siempre satisface la regla del producto ni la regla de la cadena, puesto que puede que el producto o la composición de funciones en $L^1_\text{loc}(\left]a,b\right[)$ no esté en $L^1_\text{loc}(\left]a,b\right[)$. Sin embargo, en situaciones buenas sí que se cumplen las mismas reglas.
\end{observacion}

\begin{ejemplo}
    Vamos a calcular la derivada débil de la función $f:\left]-1,1\right[\to \mathbb{R}$ valor absoluto, $f(x) = |x|$. Guiándonos por su definición, para $\varphi \in C_c^1(\left]-1,1\right[)$ escribimos:
    \begin{align*}
        \int_{-1}^{1} |x|\varphi'(x)~dx  &= -\int_{-1}^{0} x\varphi'(x)~dx + \int_{0}^{1} x\varphi(x)~dx  \\
                                         &\AstIg \int_{-1}^{0} \varphi(x)~dx - \left[x\varphi(x)\right]_{-1}^0 - \int_{0}^{1} \varphi(x)~dx + [x\varphi(x)]_0^1 \\
                                         &= \int_{-1}^{0} \varphi(x)~dx -\int_{0}^{1} \varphi(x)~dx  -\cancelto{0}{\varphi(-1)} + \cancelto{0}{\varphi(1)} = -\int_{-1}^{1} sgn(x)\varphi(x)~dx 
    \end{align*}
    donde en $(\ast)$ hemos usado integración por partes, algo que no podíamos haber aplicado a $|x|\varphi'(x)$ por no ser derivable en $0$, pero sí que lo podemos aplicar en los intervalos $\left]-1,0\right[$ y $\left]0,1\right[$, como lo hemos hecho. Hemos considerado:
    \begin{equation*}
        sgn(x) = \left\{\begin{array}{ll}
            -1 & \text{si\ } x<0 \\
            0 & \text{si\ } x=0 \\
            1 & \text{si\ } x>0
        \end{array}\right. 
    \end{equation*}
    Aunque en realidad podríamos haber considerado cualquier función igual a esta casi por doquier.\\

    \noindent
    Si hubiéramos sido algo más espabilados, podríamos haber dicho que $f\big|_{\left]-1,0\right[}$ y $f\big|_{\left]0,1\right[}$ son derivables y conocemos sus derivadas, con lo que podríamos haber obtenido la función signo (o cualquiera igual casi por doquier) como derivada débil directamente.
\end{ejemplo}

\noindent
A partir del ejemplo anterior, no es difícil darse cuenta de que cualquier función con un número de ``picos'' contenido en un conjunto de medida nula es débilmente derivable.

\begin{ejemplo}
    Tratemos de calcular la derivada débil de $sgn(x)$ en $\left]-1,1\right[$ según la definición hecha en el ejemplo anterior:
    \begin{align*}
        \int_{-1}^{1} sgn(x)\varphi'(x)~dx &= \int_{-1}^{0} -\varphi'(x)~dx + \int_{0}^{1} \varphi'(x)~dx \\
                                           &= \left[-\varphi(0) + \cancelto{0}{\varphi(-1)}\right] + \left[\cancelto{0}{\varphi(1)} - \varphi(0)\right] = -2\varphi(0) \quad \forall \varphi \in C_c^1(\left]-1,1\right[)
    \end{align*}
    Ahora, buscamos $g$ para tener:
    \begin{equation*}
        -2\varphi(0) = -\int_{-1}^{1} g(x)\varphi(x)~dx 
    \end{equation*}
    Sin embargo, no existe ninguna función $g\in L^1_\text{loc}(\left]-1,1\right[)$ que lo verifique, por lo que $sgn(x)$ no es débilmente derivable en $\left]-1,1\right[$.\\

    \noindent
    Veamos que no existe dicha $g\in L^1_\text{loc}(\left]a,b\right[)$, por reducción al absurdo, supongamos que existe $g\in L^1_\text{loc}(\left]a,b\right[)$ con:
    \begin{equation*}
        2\varphi(0) = \int_{-1}^{1} g(x)\varphi(x)~dx  \qquad \forall \varphi \in C_c^1(\left]-1,1\right[)
    \end{equation*}
    Así, tenemos que:
    \begin{align*}
        \int_{-1}^{0} g(x)\varphi(x)~dx &= \int_{-1}^{1} g(x)\varphi(x)~dx = -2\varphi(0) = 0 \qquad \forall \varphi \in C_c^1(\left]-1,0\right[) \\
        \int_{0}^{1} g(x)\varphi(x)~dx &= \int_{-1}^{1} g(x)\varphi(x)~dx = -2\varphi(0) = 0 \qquad \forall \varphi \in C_c^1(\left]0,1\right[)
    \end{align*}
    Usando el Lema Fundamental del Cálculo de Variaciones, de la primera línea concluimos que $g\equiv 0$ casi por doquier en $\left]0,1\right[$ y de la segunda que $g\equiv 0$ casi por doquier en $\left]-1,0\right[$, de donde $g\equiv 0$ casi por doquier en $\left]-1,1\right[$.\\

    \noindent
    Si tomamos ahora $\varphi \in C_c^1(\left]a,b\right[)$ con $\varphi(0) = 1$ obtenemos que:
    \begin{equation*}
        0 = \int_{-1}^{1} g(x)\varphi(x)~dx  = 2\varphi(0) = 2
    \end{equation*}
    Hemos llegado a \underline{contradicción}.
\end{ejemplo}

\noindent
Introducimos ahora la función de Heaviside (o ``función escalón'', \textit{step function}) en $L^1_\text{loc}(\left]-1,1\right[)$:
\begin{equation*}
    H(x) = \left\{\begin{array}{ll}
        1 & \text{si\ } x>0 \\
        0 & \text{si\ } x<0
    \end{array}\right. 
\end{equation*}
Obtendremos que:
\begin{equation*}
    \int_{-1}^{1} H(x)\varphi'(x)~dx = -\varphi(0) \qquad \forall \varphi \in C_c^1(\left]-1,1\right[)
\end{equation*}
Por lo que de manera análoga al ejemplo anterior, llegaremos a que $H$ no es débilmente derivable.\\

\noindent
Sin embargo, existe una entidad matemática (no es una función), la delta de Dirac, $\delta(x)$, tal que:
\begin{equation*}
    \int_{-1}^{1} \varphi(x)\delta(x)~dx = \varphi(0) \qquad \varphi \in C(\left]a,b\right[)
\end{equation*}
Para esto necesitamos:
\begin{itemize}
    \item Bien interpretar que $\delta$ es un funcional lineal sobre $\varphi$.
    \item Bien cambiar la integral de Lebesgue por una integral con otra medida.
\end{itemize}
